\documentclass[11pt,letterpaper]{article}

% --- Paquetes de formato y tipografía ---
\usepackage[spanish,es-tabla]{babel}
\usepackage[utf8]{inputenc}
\usepackage[T1]{fontenc}
\usepackage{geometry}
\geometry{top=2.5cm, bottom=2.5cm, left=2.5cm, right=2.5cm}
\usepackage{amsmath,amssymb,amsfonts}
\usepackage{graphicx}
\graphicspath{{doc/}{./}}
\usepackage{booktabs}
\usepackage{tabularx}
\usepackage{listings}
\usepackage{xcolor}
\usepackage{hyperref}
\usepackage{tikz}
\usetikzlibrary{shapes.geometric, arrows, positioning, fit, calc}

% --- Configuración de enlaces ---
\hypersetup{
    colorlinks=true,
    linkcolor=blue!70!black,
    citecolor=blue!70!black,
    urlcolor=blue!70!black
}

% --- Configuración para listados de código y YAML ---
\definecolor{codegreen}{rgb}{0,0.6,0}
\definecolor{codegray}{rgb}{0.5,0.5,0.5}
\definecolor{codepurple}{rgb}{0.58,0,0.82}
\definecolor{backcolour}{rgb}{0.97,0.97,0.97}

\lstdefinestyle{mystyle}{
    backgroundcolor=\color{backcolour},   
    commentstyle=\color{codegreen},
    keywordstyle=\color{magenta},
    numberstyle=\tiny\color{codegray},
    stringstyle=\color{codepurple},
    basicstyle=\ttfamily\footnotesize,
    breakatwhitespace=false,         
    breaklines=true,                 
    captionpos=b,                    
    keepspaces=true,                 
    numbers=left,                    
    numbersep=5pt,                  
    showspaces=false,                
    showstringspaces=false,
    showtabs=false,                  
    tabsize=2,
    frame=single
}
\lstset{style=mystyle}

% --- Estilos de diagrama TikZ ---
\tikzstyle{startstop} = [rectangle, rounded corners, minimum width=2.5cm, minimum height=0.8cm, text centered, draw=black, fill=red!20]
\tikzstyle{process} = [rectangle, minimum width=2.8cm, minimum height=0.8cm, text centered, draw=black, fill=orange!20]
\tikzstyle{decision} = [diamond, minimum width=2.5cm, minimum height=0.8cm, text centered, draw=black, fill=green!20, aspect=2]
\tikzstyle{io} = [trapezium, trapezium left angle=70, trapezium right angle=110, minimum width=2.5cm, minimum height=0.8cm, text centered, draw=black, fill=blue!20]
\tikzstyle{arrow} = [thick,->,>=stealth]

% ========================================================
% Portada institucional (estilo UVG)
% ========================================================
\begin{document}

\begin{titlepage}
    \centering
    \vspace*{0.3cm}
    
    {\bfseries\Large UNIVERSIDAD DEL VALLE DE GUATEMALA\par}
    \vspace{0.25cm}
    {\large Facultad de Ingeniería\par}
    {\large Departamento de Ciencias de la Computación\par}
    
    \vspace{0.8cm}
    \includegraphics[width=4.5cm]{logo_uvg.png}\par
    \vspace{0.8cm}
    
    {\large CC2019 -- Teoría de la Computación\par}
    {\large Ciclo 2 de 2026\par}
    
    \vspace{1.2cm}
    
    {\bfseries\LARGE Proyecto 3: Simulador universal de\\ máquinas de Turing\par}
    
    \vspace{0.6cm}
    
    {\large Implementación formal bajo el modelo de cuádruplas, programación defensiva y análisis asintótico de complejidad temporal\par}
    
    \vfill
    
    {\bfseries\large Integrantes del grupo:\par}
    \vspace{0.3cm}
    {\large
    Samuel Mejía -- 23442\par
    % Descomentar y completar si el grupo tiene 2 o 3 integrantes:
    % Nombre Integrante 2 -- Carnet\par
    % Nombre Integrante 3 -- Carnet\par
    }
    
    \vspace{1.2cm}
    
    {\small
    \textbf{Repositorio:} \url{https://github.com/SamMejia23442/proyectofinal-teoria-compu.git}\par
    \vspace{0.15cm}
    \textbf{Presentación:} \url{https://www.youtube.com/watch?v=VIDEO_ID_AQUI}\par
    \vspace{0.25cm}
    Guatemala, noviembre de 2026\par
    }
    
\end{titlepage}

\newpage
\thispagestyle{plain}

\begin{abstract}
En este informe se presenta el diseño, especificación formal e implementación de un simulador universal de máquinas de Turing deterministas estándar de una cinta infinita bidireccional y cabezal único. El motor interpreta especificaciones en formato YAML basadas en la convención formal de cuádruplas $(\text{origen}, \text{símbolo}, \text{acción}, \text{destino})$, incorporando programación defensiva integral contra especificaciones mal formadas, transiciones no deterministas o entradas inválidas. Asimismo, se documenta la computación paso a paso mediante descripciones instantáneas (IDs), el control de ejecución preventivo con cancelación de ciclos indefinidos, y la demostración formal de la tasa de crecimiento asintótico para dos máquinas originales: un \textit{predecesor binario} ($O(n)$) y un \textit{verificador de palíndromos} ($O(n^2)$).
\end{abstract}

\tableofcontents
\newpage

% ========================================================
\section{Introducción y marco teórico}
% ========================================================

La máquina de Turing (MT), formulada por Alan Turing en 1936, constituye el modelo matemático fundacional de la computación universal. Formalmente, una máquina de Turing determinista estándar se define como una 7-tupla:
\begin{equation}
M = (Q, \Sigma, \Gamma, \delta, q_0, B, F)
\end{equation}
donde $Q$ representa el conjunto finito de estados, $\Sigma$ es el alfabeto de entrada, $\Gamma$ es el alfabeto de la cinta ($\Sigma \subset \Gamma$), $B \in \Gamma \setminus \Sigma$ es el símbolo de espacio en blanco, $q_0 \in Q$ es el estado de inicio y $F \subseteq Q$ es el conjunto de estados de aceptación o parada.

\subsection{El modelo de cuádruplas de Post-Davis}
En los cursos modernos de Teoría de la Computación, la función de transición tradicional basada en quíntuplas $\delta: Q \times \Gamma \to Q \times \Gamma \times \{L, R\}$ se descompone con frecuencia en el \textbf{modelo formal de cuádruplas} (introducido por Emil Post y Martin Davis). En este modelo, una transición se expresa como:
\begin{equation}
(q_i, s_j, a_k, q_l)
\end{equation}
donde la acción $a_k$ es estrictamente atómica:
\begin{enumerate}
    \item \textbf{Movimiento del cabezal:} $a_k \in \{L, R\}$. El cabezal se desplaza una casilla a la izquierda ($L$) o a la derecha ($R$), preservando el símbolo actual de la cinta.
    \item \textbf{Escritura sobre la cinta:} $a_k \in \Gamma$. El cabezal sobreescribe la casilla actual con el símbolo $a_k$, permaneciendo en la misma posición de la cinta.
\end{enumerate}

Ambos modelos son computacionalmente equivalentes: cualquier quíntupla $(q_i, s_j, s'_j, D, q_l)$ se simula de manera idéntica mediante dos cuádruplas sucesivas incorporando un estado intermedio efímero:
\begin{equation}
(q_i, s_j, s'_j, q_{\text{temp}}) \quad \text{y} \quad (q_{\text{temp}}, s'_j, D, q_l)
\end{equation}

\subsection{Descripciones instantáneas (IDs)}
Para registrar formalmente el estado de la máquina en cada instante temporal de la computación, se emplea la notación de \textbf{descripciones instantáneas (IDs)}. Una ID se denota por:
\begin{equation}
\alpha [q] \beta
\end{equation}
donde $\alpha \in \Gamma^*$ es la secuencia de símbolos no blancos a la izquierda del cabezal, $q \in Q$ es el estado actual de la unidad de control, y $\beta \in \Gamma^+$ representa la porción de la cinta que comienza exactamente en la casilla actual leída por el cabezal y concluye en el último símbolo no blanco a la derecha.

% ========================================================
\section{Especificación de sintaxis YAML y símbolo de blanco}
% ========================================================

\subsection{Sintaxis formal del archivo YAML}
La especificación de la máquina de Turing se provee mediante un archivo con sintaxis YAML bajo la clave raíz obligatoria \texttt{mt}. La estructura esperada se define a continuación:

\begin{lstlisting}[language=SQL, caption={Plantilla formal de especificación YAML.}]
mt:
  estados: [q0, q1, q2, q_fin]          # Lista exhaustiva de estados
  alfabeto: ["0", "1", B, Y, N]         # Alfabeto unificado (entrada y cinta)
  estado_inicial: q0                    # Estado de arranque de la maquina
  blanco: B                             # Simbolo para espacios vacios
  transiciones:                         # Lista de cuadruplas separadas por espacios
    - q0 0 R q0
    - q0 1 R q0
    - q0 B L q1
    - q1 1 0 q2
\end{lstlisting}

\subsection{Identificación del símbolo para espacios en blanco}
\begin{itemize}
    \item Por estándar del proyecto, el símbolo que representa los espacios en blanco sobre la cinta es la letra mayúscula \textbf{\texttt{B}} (\textit{Blank}).
    \item El simulador admite de forma opcional la clave explícita \texttt{blanco} en el archivo YAML para definir símbolos alternativos (como el guion bajo \texttt{\_}). Si la clave se omite, el analizador detecta automáticamente \texttt{B} o \texttt{\_} dentro de la lista de alfabeto declarada.
    \item \textbf{Cinta infinita:} Las posiciones no inicializadas de la cinta hacia la izquierda ($-\infty, \dots, -1$) y hacia la derecha ($n, \dots, +\infty$) se consideran implícitamente pobladas por el símbolo de blanco \texttt{B}.
\end{itemize}

\subsection{Mecanismos de programación defensiva}
El simulador incorpora un módulo de validación estricta (\texttt{parser.py}) que rechaza cualquier especificación inconsistente antes de iniciar la ejecución:
\begin{table}[h!]
\centering
\small
\begin{tabularx}{\textwidth}{lX}
\toprule
\textbf{Condición defensiva} & \textbf{Validación y acción del simulador} \\
\midrule
Estructura y sintaxis YAML & Verifica existencia del archivo, formato sintáctico y presencia de las claves obligatorias (\texttt{estados}, \texttt{alfabeto}, \texttt{estado\_inicial}, \texttt{transiciones}). \\
Integridad de conjuntos & Valida que los estados y los caracteres del alfabeto sean únicos, no nulos y tengan longitud 1 en el alfabeto. \\
Estado inicial & Asegura que \texttt{estado\_inicial} pertenezca formalmente al conjunto \texttt{estados}. \\
Formato de cuádruplas & Comprueba que cada regla posea exactamente 4 elementos separados por espacios: \texttt{origen simbolo accion destino}. \\
Consistencia de transición & Corrobora que \texttt{origen} y \texttt{destino} $\in Q$, que \texttt{simbolo} $\in \Gamma$ y que la acción sea válida ($R, L$ o un símbolo $\in \Gamma$). \\
No determinismo & Detecta reglas duplicadas para un mismo par $(q_{\text{actual}}, s_{\text{leído}})$ y aborta inmediatamente con mensaje de error específico. \\
Validación de la entrada & Contrasta cada carácter de la cadena recibida contra el alfabeto declarado antes de inicializar la cinta. \\
\bottomrule
\end{tabularx}
\caption{Matriz de validaciones y programación defensiva implementada.}
\end{table}

% ========================================================
\section{Manual de uso del simulador}
% ========================================================

\subsection{Requisitos y configuración del entorno}
El software está programado en Python 3 (versión 3.8 o superior) y requiere la instalación del paquete \texttt{PyYAML}:
\begin{lstlisting}[language=bash]
# Clonar el repositorio
git clone https://github.com/SamMejia23442/proyectofinal-teoria-compu.git
cd proyectofinal-teoria-compu

# Instalar dependencias
pip install -r requirements.txt
\end{lstlisting}

\subsection{Modo interactivo (por menú)}
Al invocar el programa sin parámetros en la terminal, se despliega una consola interactiva:
\begin{lstlisting}[language=bash]
python3 main.py
\end{lstlisting}
El menú ofrece la selección directa de las máquinas de ejemplo, ingreso de cadenas de prueba personalizadas y opción para habilitar la simulación paso a paso.

\subsection{Modo línea de comandos (CLI)}
Para automatización, pruebas unitarias o ejecución directa:
\begin{lstlisting}[language=bash]
# Sintaxis general:
python3 main.py <archivo_yaml> [cadena_entrada] [opciones]

# Ejemplo 1: Predecesor binario para la entrada 1100
python3 main.py ejemplos/predecesor_binario.yaml 1100

# Ejemplo 2: Verificador de palíndromos para 10101
python3 main.py ejemplos/verificador_palindromos.yaml 10101

# Ejemplo 3: Ejecución paso a paso interactiva
python3 main.py ejemplos/verificador_palindromos.yaml 10101 -s
\end{lstlisting}

\subsubsection{Banderas y parámetros disponibles}
\begin{itemize}
    \item \texttt{-s}, \texttt{--paso-a-paso}: Pausa la ejecución tras cada transición. Permite avanzar manualmente presionando \texttt{Enter}, pasar a modo continuo automático ingresando \texttt{t}, o cancelar la computación con \texttt{c}.
    \item \texttt{--sin-visual}: Omite la tira gráfica de casillas de la cinta en consola, imprimiendo únicamente la lista formal de IDs.
    \item \texttt{--limite <N>}: Configura un límite preventivo de $N$ pasos para interrumpir bucles infinitos en máquinas de Turing generales.
    \item \texttt{--no-color}: Desactiva los colores ANSI para terminales planas o redirección a archivos de texto.
\end{itemize}

\subsection{Control de ejecución y cancelación}
De acuerdo con las restricciones del proyecto, el simulador no impone restricciones artificiales a la cinta ni a la cantidad de pasos. No obstante, para casos con resultados indefinidos (lenguajes indecidibles o bucles infinitos), el usuario puede presionar \textbf{\texttt{Ctrl+C}} en cualquier instante. El simulador intercepta la señal \texttt{SIGINT} de forma limpia, reporta el motivo de parada como cancelación por el usuario, e imprime el resumen de los pasos ejecutados y la ID alcanzada sin generar volcados de pila (\textit{stack traces}) descontrolados.

% ========================================================
\section{Diagrama de flujo del algoritmo de simulación}
% ========================================================

La Figura~\ref{fig:flowchart} esquematiza el ciclo de vida computacional del simulador, desde la lectura de la especificación YAML hasta la extracción del resultado final en la cinta.

\begin{figure}[h!]
\centering
\begin{tikzpicture}[node distance=1.3cm, scale=0.88, every node/.style={scale=0.88}]
    \node (start) [startstop] {Inicio};
    \node (read_yaml) [process, below of=start] {Cargar archivo YAML};
    \node (val_yaml) [decision, below of=read_yaml, yshift=-0.4cm] {¿YAML válido?\\Defensa OK?};
    \node (err_spec) [startstop, right of=val_yaml, xshift=2.5cm, fill=red!40] {Abortar: Error\\de especificación};
    
    \node (init_tape) [process, below of=val_yaml, yshift=-0.5cm] {Inicializar cinta infinita\\$w \in \Gamma^*$, Cabezal en pos $0$};
    \node (val_input) [decision, below of=init_tape, yshift=-0.4cm] {¿Símbolos de $w$\\pertenecen a $\Sigma$?};
    \node (err_input) [startstop, right of=val_input, xshift=2.5cm, fill=red!40] {Abortar: Entrada\\inválida};
    
    \node (print_id) [io, below of=val_input, yshift=-0.5cm] {Registrar e imprimir ID\\$\alpha [q] \beta$ y cinta};
    \node (check_trans) [decision, below of=print_id, yshift=-0.6cm] {¿Existe $\delta(q, s)$?};
    
    \node (check_type) [decision, below of=check_trans, yshift=-0.7cm] {¿Acción es\\movimiento o escritura?};
    \node (do_move) [process, left of=check_type, xshift=-2.2cm] {Desplazar cabezal\\$pos \leftarrow pos \pm 1$};
    \node (do_write) [process, right of=check_type, xshift=2.2cm] {Sobreescribir celda\\$cinta[pos] \leftarrow a_k$};
    
    \node (update_state) [process, below of=check_type, yshift=-0.7cm] {$q \leftarrow q_{\text{destino}}$, Pasos $\leftarrow$ Pasos $+ 1$};
    \node (check_cancel) [decision, below of=update_state, yshift=-0.5cm] {¿Cancelación\\o Ctrl+C?};
    
    \node (halt) [process, right of=check_trans, xshift=4.2cm] {Parada de la máquina\\(Halt)};
    \node (extract) [io, below of=halt, yshift=-1.0cm] {Extraer resultado\\limpiando blancos $B$};
    \node (end) [startstop, below of=extract] {Fin de simulación};

    % Conexiones
    \draw [arrow] (start) -- (read_yaml);
    \draw [arrow] (read_yaml) -- (val_yaml);
    \draw [arrow] (val_yaml) -- node[anchor=south] {No} (err_spec);
    \draw [arrow] (val_yaml) -- node[anchor=east] {Sí} (init_tape);
    \draw [arrow] (init_tape) -- (val_input);
    \draw [arrow] (val_input) -- node[anchor=south] {No} (err_input);
    \draw [arrow] (val_input) -- node[anchor=east] {Sí} (print_id);
    
    \draw [arrow] (print_id) -- (check_trans);
    \draw [arrow] (check_trans) -- node[anchor=south] {No} (halt);
    \draw [arrow] (check_trans) -- node[anchor=east] {Sí} (check_type);
    
    \draw [arrow] (check_type) -- node[anchor=south] {Movimiento} (do_move);
    \draw [arrow] (check_type) -- node[anchor=south] {Escritura} (do_write);
    
    \draw [arrow] (do_move) |- (update_state);
    \draw [arrow] (do_write) |- (update_state);
    
    \draw [arrow] (update_state) -- (check_cancel);
    \draw [arrow] (check_cancel) -- node[anchor=south] {Sí} ++(5.5,0) |- (extract);
    \draw [arrow] (check_cancel) -- node[anchor=east] {No} ++(-4.5,0) |- (print_id);
    
    \draw [arrow] (halt) -- (extract);
    \draw [arrow] (extract) -- (end);
\end{tikzpicture}
\caption{Diagrama de flujo de la ejecución y control del simulador de máquinas de Turing.}
\label{fig:flowchart}
\end{figure}

% ========================================================
\section{Máquinas de Turing de ejemplo y análisis asintótico}
% ========================================================

En estricto cumplimiento con la directiva de no repetir algoritmos con otros grupos de trabajo, se seleccionaron, diseñaron e implementaron dos máquinas de Turing originales:
\begin{enumerate}
    \item \textbf{Máquina 1 (aritmética):} Predecesor binario ($n - 1$).
    \item \textbf{Máquina 2 (reconocimiento de lenguajes):} Verificador de palíndromos sobre $\{0, 1\}^*$.
\end{enumerate}

% --------------------------------------------------------
\subsection{Máquina 1: Predecesor binario ($n - 1$)}
% --------------------------------------------------------

\subsubsection{Definición formal y propósito}
Dado un número binario positivo $N \in \{0, 1\}^+$ escrito en la cinta, la máquina resta 1 a su valor aritmético ($N - 1$) propagando los préstamos (\textit{borrows}) necesarios de derecha a izquierda, y deja el resultado exacto sobre la cinta con el cabezal posicionado en el inicio del número.

Formalmente:
\begin{itemize}
    \item $Q = \{q_0, q_1, q_{\text{borrow}}, q_2, q_{\text{fin}}\}$
    \item $\Sigma = \{0, 1\}$, $\Gamma = \{0, 1, B\}$, $q_0 = q_0$, $B = B$
    \item Parada normal: Al alcanzar el estado $q_{\text{fin}}$ sobre el primer bit resultante.
\end{itemize}

\subsubsection{Tabla de transiciones (cuádruplas)}
\begin{table}[h!]
\centering
\small
\begin{tabular}{ccccc}
\toprule
\textbf{Cuádrupla} & \textbf{Estado origen} & \textbf{Símbolo leído} & \textbf{Acción} & \textbf{Estado siguiente} \\
\midrule
1 & $q_0$ & \texttt{0} & $R$ (mover a la derecha) & $q_0$ \\
2 & $q_0$ & \texttt{1} & $R$ (mover a la derecha) & $q_0$ \\
3 & $q_0$ & \texttt{B} & $L$ (mover a la izquierda) & $q_1$ \\
\midrule
4 & $q_1$ & \texttt{1} & \texttt{0} (escribir 0) & $q_2$ \\
5 & $q_1$ & \texttt{0} & \texttt{1} (escribir 1) & $q_{\text{borrow}}$ \\
6 & $q_{\text{borrow}}$ & \texttt{1} & $L$ (mover a la izquierda) & $q_1$ \\
\midrule
7 & $q_2$ & \texttt{0} & $L$ (mover a la izquierda) & $q_2$ \\
8 & $q_2$ & \texttt{1} & $L$ (mover a la izquierda) & $q_2$ \\
9 & $q_2$ & \texttt{B} & $R$ (mover a la derecha) & $q_{\text{fin}}$ \\
\bottomrule
\end{tabular}
\caption{Función de transición para la máquina predecesor binario.}
\end{table}

\subsubsection{Trazas de prueba ejecutadas}
\begin{itemize}
    \item \textbf{Prueba 1 ($1100_2 \to 1011_2$):} 13 pasos de computación.
    \begin{quote}\footnotesize
    $[q_0]1100 \vdash 1[q_0]100 \vdash 11[q_0]00 \vdash 110[q_0]0 \vdash 1100[q_0]B \vdash 110[q_1]0 \vdash 110[q_{\text{borrow}}]1 \vdash 11[q_1]01 \vdash 11[q_{\text{borrow}}]11 \vdash 1[q_1]111 \vdash 1[q_2]011 \vdash [q_2]1011 \vdash [q_2]B1011 \vdash [q_{\text{fin}}]1011$.
    \end{quote}
    \item \textbf{Prueba 2 ($1010_2 \to 1001_2$):} 12 pasos de computación.
    \item \textbf{Prueba 3 ($1000_2 \to 0111_2$):} 14 pasos de computación.
    \item \textbf{Prueba 4 ($1_2 \to 0_2$):} 5 pasos de computación.
\end{itemize}

\subsubsection{Análisis de complejidad asintótica}
Sea $n = |w|$ la longitud en bits de la cadena de entrada binaria. El algoritmo consta de tres fases secuenciales bien delimitadas:
\begin{enumerate}
    \item \textbf{Fase 1 (búsqueda del extremo derecho):} El cabezal inicia en el bit más significativo ($pos = 0$) y avanza a la derecha hasta leer el primer espacio en blanco en $pos = n$. Se ejecutan exactamente $n$ movimientos $R$ en $q_0$ y $1$ movimiento $L$ para ubicarse en el bit menos significativo ($pos = n-1$). Esta fase requiere exactamente $n + 1$ pasos.
    \item \textbf{Fase 2 (resta y propagación de préstamos):} Si el número termina en $k$ ceros ($0 \le k < n$), cada cero requiere 1 escritura (cambia a \texttt{1}) y 1 movimiento $L$ hacia la izquierda, tomando $2k$ pasos. Al alcanzar el primer bit \texttt{1}, se realiza 1 escritura (cambia a \texttt{0}). Dado que a lo sumo se recorren todos los $n$ bits de la entrada, esta fase ejecuta a lo sumo $2(n-1) + 1 = 2n - 1$ pasos.
    \item \textbf{Fase 3 (retorno del cabezal al inicio):} El cabezal se encuentra a lo sumo en la posición $pos = n - k - 1$ y retrocede hacia la izquierda hasta encontrar el blanco en $pos = -1$, seguido de un movimiento $R$ para reposicionarse en $pos = 0$. Esto requiere a lo sumo $n + 1$ pasos.
\end{enumerate}

Sumando las tres fases, el tiempo total de ejecución $T(n)$ en el peor de los casos satisface:
\begin{equation}
T_{\text{peor}}(n) = (n + 1) + (2n - 1) + (n + 1) = 4n + 1
\end{equation}
Para cualquier entrada de tamaño $n \ge 1$:
\begin{equation}
c_1 n \le T(n) \le c_2 n
\end{equation}
Por consiguiente, la tasa de crecimiento asintótico del predecesor binario es estrictamente:
\begin{equation}
\mathbf{\Theta(n) \quad \text{o en cota superior} \quad O(n)}
\end{equation}

% --------------------------------------------------------
\subsection{Máquina 2: Verificador de palíndromos sobre $\{0, 1\}$}
% --------------------------------------------------------

\subsubsection{Definición formal y propósito}
Determina si una cadena binaria $w \in \{0, 1\}^*$ es simétrica ($w = w^R$). La máquina compara el primer símbolo con el último; si coinciden, los consume (sustituyéndolos por espacios en blanco $B$) y repite el proceso hacia el centro. Si todos los pares coinciden, escribe \textbf{\texttt{Y}} (\textit{Yes / Aceptada}). Si ocurre una discrepancia, escribe \textbf{\texttt{N}} (\textit{No / Rechazada}) y limpia el residuo de la cinta.

Formalmente:
\begin{itemize}
    \item $Q = \{q_{\text{start}}, q_{\text{saw0\_move}}, q_{\text{saw1\_move}}, q_{\text{find\_end0}}, q_{\text{find\_end1}}, q_{\text{check0}}, q_{\text{check1}}, q_{\text{return\_move}}, q_{\text{return\_left}}, q_{\text{clean\_start}}, q_{\text{clean\_left}}, q_{\text{clean\_move}}, q_{\text{halt}}\}$
    \item $\Sigma = \{0, 1\}$, $\Gamma = \{0, 1, B, Y, N\}$, $q_0 = q_{\text{start}}$, $B = B$
    \item Parada: Estado $q_{\text{halt}}$, con el veredicto (\texttt{Y} o \texttt{N}) como único contenido en la cinta.
\end{itemize}

\subsubsection{Tabla de transiciones (cuádruplas)}
\begin{table}[h!]
\centering
\footnotesize
\begin{tabular}{ccccc}
\toprule
\textbf{\#} & \textbf{Estado origen} & \textbf{Símbolo leído} & \textbf{Acción} & \textbf{Estado siguiente} \\
\midrule
1 & $q_{\text{start}}$ & \texttt{0} & \texttt{B} & $q_{\text{saw0\_move}}$ \\
2 & $q_{\text{start}}$ & \texttt{1} & \texttt{B} & $q_{\text{saw1\_move}}$ \\
3 & $q_{\text{start}}$ & \texttt{B} & \texttt{Y} & $q_{\text{halt}}$ \\
\midrule
4 & $q_{\text{saw0\_move}}$ & \texttt{B} & $R$ & $q_{\text{find\_end0}}$ \\
5 & $q_{\text{saw1\_move}}$ & \texttt{B} & $R$ & $q_{\text{find\_end1}}$ \\
\midrule
6 & $q_{\text{find\_end0}}$ & \texttt{0} / \texttt{1} & $R$ & $q_{\text{find\_end0}}$ \\
7 & $q_{\text{find\_end0}}$ & \texttt{B} & $L$ & $q_{\text{check0}}$ \\
8 & $q_{\text{find\_end1}}$ & \texttt{0} / \texttt{1} & $R$ & $q_{\text{find\_end1}}$ \\
9 & $q_{\text{find\_end1}}$ & \texttt{B} & $L$ & $q_{\text{check1}}$ \\
\midrule
10 & $q_{\text{check0}}$ & \texttt{0} & \texttt{B} & $q_{\text{return\_move}}$ \\
11 & $q_{\text{check0}}$ & \texttt{B} & \texttt{Y} & $q_{\text{halt}}$ \\
12 & $q_{\text{check0}}$ & \texttt{1} & \texttt{N} & $q_{\text{clean\_start}}$ \\
13 & $q_{\text{check1}}$ & \texttt{1} & \texttt{B} & $q_{\text{return\_move}}$ \\
14 & $q_{\text{check1}}$ & \texttt{B} & \texttt{Y} & $q_{\text{halt}}$ \\
15 & $q_{\text{check1}}$ & \texttt{0} & \texttt{N} & $q_{\text{clean\_start}}$ \\
\midrule
16 & $q_{\text{return\_move}}$ & \texttt{B} & $L$ & $q_{\text{return\_left}}$ \\
17 & $q_{\text{return\_left}}$ & \texttt{0} / \texttt{1} & $L$ & $q_{\text{return\_left}}$ \\
18 & $q_{\text{return\_left}}$ & \texttt{B} & $R$ & $q_{\text{start}}$ \\
\midrule
19 & $q_{\text{clean\_start}}$ & \texttt{N} & $L$ & $q_{\text{clean\_left}}$ \\
20 & $q_{\text{clean\_left}}$ & \texttt{0} / \texttt{1} & \texttt{B} & $q_{\text{clean\_move}}$ \\
21 & $q_{\text{clean\_left}}$ & \texttt{B} & $R$ & $q_{\text{halt}}$ \\
22 & $q_{\text{clean\_move}}$ & \texttt{B} & $L$ & $q_{\text{clean\_left}}$ \\
\bottomrule
\end{tabular}
\caption{Función de transición resumida para el verificador de palíndromos.}
\end{table}

\subsubsection{Trazas de prueba ejecutadas}
\begin{itemize}
    \item \textbf{Palíndromo impar ($10101 \to \texttt{Y}$):} 26 pasos de simulación. Coinciden extremos en \texttt{1}, luego en \texttt{0}, y el símbolo central \texttt{1} detecta blanco adyacente y escribe \texttt{Y}.
    \item \textbf{Palíndromo par ($1001 \to \texttt{Y}$):} 19 pasos de simulación. Ambos pares coinciden y la cinta vacía final transiciona a \texttt{Y}.
    \item \textbf{No palíndromo ($10110 \to \texttt{N}$):} 16 pasos de simulación. El extremo izquierdo (\texttt{1}) discrepa con el derecho (\texttt{0}); escribe \texttt{N}, limpia los bits internos restantes y concluye.
    \item \textbf{Longitudes unitarias y cadena vacía ($1 \to \texttt{Y}$, $0 \to \texttt{Y}$, $\epsilon \to \texttt{Y}$):} 4, 4 y 1 pasos respectivamente.
\end{itemize}

\subsubsection{Análisis de complejidad asintótica}
Sea $n = |w|$ la longitud de la cadena de entrada. En el peor de los casos (la cadena es efectivamente un palíndromo), el algoritmo realiza $\lceil n/2 \rceil$ iteraciones de comparación. En la iteración $i$-ésima ($i = 0, 1, \dots, \lceil n/2 \rceil - 1$):
\begin{enumerate}
    \item Borra el símbolo izquierdo: 1 paso.
    \item Recorre los símbolos restantes hasta el extremo derecho: la porción restante tiene longitud $n - 2i - 1$. Requiere $n - 2i$ desplazamientos a la derecha.
    \item Compara y borra el símbolo derecho coincidente: 1 paso de lectura/escritura.
    \item Retorna hacia la izquierda a través de los $n - 2i - 2$ símbolos intermedios hasta encontrar el blanco izquierdo: $n - 2i - 1$ desplazamientos a la izquierda.
\end{enumerate}

Por tanto, en cada ronda se ejecutan aproximadamente:
\begin{equation}
T_{\text{ronda}}(i) = 2(n - 2i) + c
\end{equation}
Sumando sobre todas las $\approx \frac{n}{2}$ iteraciones:
\begin{equation}
T(n) = \sum_{i=0}^{n/2} [2(n - 2i) + c] = 2n \left(\frac{n}{2}\right) - 4 \sum_{i=0}^{n/2} i + c \frac{n}{2} = n^2 - 4 \frac{(n/2)^2}{2} + O(n) = \frac{1}{2} n^2 + O(n)
\end{equation}
Por ende, la tasa de crecimiento asintótico para el verificador de palíndromos en una máquina de Turing de una cinta es cuadrática:
\begin{equation}
\mathbf{\Theta(n^2) \quad \text{o en cota superior} \quad O(n^2)}
\end{equation}

% ========================================================
\section{Enlaces del proyecto y entrega}
% ========================================================

\begin{itemize}
    \item \textbf{Repositorio en GitHub:} Código fuente del simulador, pruebas unitarias, archivos de especificación YAML y documentación técnica:\\[0.2cm]
    \url{https://github.com/SamMejia23442/proyectofinal-teoria-compu.git}
    
    \item \textbf{Video de presentación en YouTube (no listado, máx. 7 minutos):}\\[0.2cm]
    \url{https://www.youtube.com/watch?v=VIDEO_ID_AQUI} \\
    \textit{(Nota: Reemplazar el identificador con el enlace del video no listado grabado por los integrantes antes de la entrega final).}
\end{itemize}

% ========================================================
\section{Conclusiones}
% ========================================================
\begin{enumerate}
    \item Se implementó con éxito un simulador determinista de máquinas de Turing de una cinta no acotada y cabezal bidireccional, capaz de interpretar formalmente funciones de transición bajo el esquema de cuádruplas $(\text{origen}, \text{símbolo}, \text{acción}, \text{destino})$.
    \item La programación defensiva implementada previene fallos en tiempo de ejecución al validar rigurosamente la consistencia de tipos, pertenencia de elementos a los alfabetos y ausencia de no determinismo antes de procesar cualquier entrada.
    \item El análisis asintótico formal concuerda con las métricas observadas durante las ejecuciones experimentales: las transformaciones que requieren un barrido simple (como el predecesor binario) operan en tiempo $O(n)$, mientras que los reconocedores de simetría en una única cinta (como el verificador de palíndromos) requieren inherentemente tiempo cuadrático $O(n^2)$ debido al costo de oscilación del cabezal entre los extremos de la cinta.
\end{enumerate}

\end{document}
